% \iffalse meta-comment
%
%% regulatory-sources.dtx
%% Copyright 2024-2026 E. Nijenhuis
%
% This work may be distributed and/or modified under the
% conditions of the LaTeX Project Public License, either version 1.3c
% of this license or (at your option) any later version.
% The latest version of this license is in
% http://www.latex-project.org/lppl.txt
% and version 1.3c or later is part of all distributions of LaTeX
% version 2005/12/01 or later.
%
% This work has the LPPL maintenance status ‘maintained’.
%
% The Current Maintainer of this work is E. Nijenhuis.
%
% This work consists of the files listed in the meta-comment of
% regulatory-struct.dtx.
%
% \fi
%
% \iffalse
%<*driver>
\ProvidesFile{regulatory-sources.dtx}
%</driver>
%<package>\NeedsTeXFormat{LaTeX2e}
%<package>\ProvidesPackage{regulatory-sources}
%<package>    [2026/09/06 0.0.5 Xerdi's Regulatory Package (External sources)]
%
%<*driver>
\documentclass[10pt,english]{ltxdoc}
%! suppress = InclusionLoop
\usepackage{regulatory}
\usepackage{tabularx}
\usepackage[english,dutch]{babel}
%% regulatory-preamble.tex
%% Copyright 2024 E. Nijenhuis
%
% This work may be distributed and/or modified under the
% conditions of the LaTeX Project Public License, either version 1.3c
% of this license or (at your option) any later version.
% The latest version of this license is in
% http://www.latex-project.org/lppl.txt
% and version 1.3c or later is part of all distributions of LaTeX
% version 2005/12/01 or later.
%
% This work has the LPPL maintenance status ‘maintained’.
%
% The Current Maintainer of this work is E. Nijenhuis.
%
% This work consists of the files regulatory.ins,
% regulatory-struct.dtx, regulatory-ref.dtx,
% regulatory-defs.dtx, regulatory-attachments.dtx,
% regulatory-md.dtx,
% regulatory-html.dtx,
% regulatory-english.def, regulatory-dutch.def,
% regulatory.tex, regulatory-preamble.tex,
% regulatory-nl.tex, regulatory-en.tex,
% example1.bib, example2.bib,
% example1-nl.tex, example2-nl.tex,
% example1-en.tex, example2-en.tex,
% md-example.tex, example.md,
% and the derived files regulatory.sty, regulatory-struct.sty,
% regulatory-defs.sty, regulatory-ref.sty,
% regulatory-attachments.sty, regulatory-md.sty
% and regulatory.4ht.
\usepackage{gitinfo-lua}
\usepackage{listings}
\usepackage{adjustbox}
\usepackage{tikz}
\usetikzlibrary{arrows.meta}

\usepackage{fontspec}
\usepackage{textcomp}

\def\projecturl{https://github.com/Xerdi/regulatory}

\definecolor{greycolor}{HTML}{BFBFBF}
\definecolor{primary}{HTML}{174A67}
\definecolor{darkcolor}{HTML}{091D28}
\definecolor{greencolor}{HTML}{176734}
\colorlet{artlabel}{primary}
\colorlet{green}{greencolor}

\lstdefinestyle{tex}{%
    language=[LaTeX]TeX,
    keywordsprefix={\\},
    alsoletter={\\},
    morekeywords={article,para,textfill,masterdocument,refdocument,loadglsdefs,printdefs,setmiddleconjunction,setlastconjunction,setrangeconjunction,setconjunction,markdownInput,glssetwidest,describe}
}

\lstdefinestyle{bash}{%
    language=bash,
    morekeywords={pdflatex,lualatex,latexmk,makeglossaries,bib2gls}
}

\lstdefinelanguage{Markdown}[LaTeX]{TeX}{%
    sensitive=false,
    morecomment=[l][\bfseries\ttfamily\color{green}]{\#},
    morecomment=[s][\bfseries\ttfamily\color{purple}]{:\ \ }{\ },
    morestring=[s][\color{black}]{\{}{\}},
    morekeywords={article,para,textfill,masterdocument,refdocument,loadglsdefs,printdefs,setmiddleconjunction,setlastconjunction,setrangeconjunction,setconjunction,markdownInput,glssetwidest,describe},
    keywordsprefix={\\},
    alsoletter={\\},
}

\lstdefinestyle{md}{%
    language=Markdown
}

\lstdefinelanguage{BibTeX}{%
    keywords={%
        article,book,collectedbook,conference,electronic,entry,ieeetranbstctl,%
        inbook,incollectedbook,incollection,injournal,inproceedings,%
        manual,mastersthesis,misc,patent,periodical,phdthesis,preamble,%
        proceedings,standard,string,techreport,unpublished%
    },
    keywordsprefix={@},
    comment=[l][\itshape]{@comment},
    sensitive=false,
}

\lstdefinestyle{bib}{%
    language=BibTeX,
    morestring=[s][\bfseries\ttfamily\color{purple}]{\ \ \ \ }{\ },
    morecomment=[n][\bfseries\ttfamily\color{green}]{\ \{}{\}},
}

\lstset{
    title=\lstname,
    basicstyle=\ttfamily,
    numberstyle=\ttfamily,
    numbersep=5pt,
    rulecolor=\color{greycolor},
    stringstyle=\color{pink},
    keywordstyle=\color{primary},
    commentstyle=\itshape\color{darkcolor},
    breaklines=true,
    captionpos=b,
    tabsize=4,
    showtabs=true
}

\newcommand\package{\texttt}
\newcommand\file{\texttt}
\newcommand\option{\texttt}

% Reads test/conformance.tex, which test/conformance.sh generates and which is
% kept in the repository: veraPDF is not available everywhere this manual is
% built, so the verdicts travel with the sources.
\newcommand\conformanceversion{?}
\newcommand\conformanceimage{?}
% A digest carries no break points of its own and does not fit a line, so this
% offers one after every character. The argument is expanded first, otherwise the
% loop sees one macro instead of the characters it stands for.
\newcommand\anywherebreakend{}
\def\anywherebreakloop#1{%
    \ifx#1\anywherebreakend\else#1\allowbreak\expandafter\anywherebreakloop\fi}
\newcommand\anywherebreak[1]{\anywherebreakloop#1\anywherebreakend}
\newcommand\conformanceimagetext{%
    \expandafter\anywherebreak\expandafter{\conformanceimage}}

% \ignorespaces, because this macro typesets nothing while its line in the
% generated file still ends in a newline, and that space would otherwise open
% the first cell of the table.
\newcommand\conformancevalidator[2]{%
    \gdef\conformanceversion{#1}\gdef\conformanceimage{#2}\ignorespaces}
\newcommand\conformancerow[6]{%
    \texttt{#1} & \texttt{#2} & \texttt{#3} & \texttt{#4} & \texttt{#5}%
    \ifthenelse{\equal{#6}{}}{}{ \footnotesize(#6)} \\}

% Points at the resulting PDF of an example, which is attached to this manual.
% Its argument is the label of the artifact the example was declared with, and it
% falls back to plain text whenever the PDF isn't there.
\newcommand\exampleresult[1]{%
    \par\noindent
    \translation{The result of this example is}{Het resultaat van dit voorbeeld is}
    \documentattachment{#1}{\file{\artifactfilename{#1}}}.\par}

\def\packagedate{\gitdate}
\def\packageversion{\gitversion}

%\def\cmd{\lstinline[style=tex]}
\def\cmditem#1{\item[\ttfamily\textbackslash{}#1]}

\newcommand\Vtextvisiblespace[1][.3em]{%
  \mbox{\kern.06em\vrule height.3ex}%
  \vbox{\hrule width#1}%
  \hbox{\vrule height.3ex}}

% Package zref-clever loads its language files lazily, which would otherwise
% happen inside a \DocInput region, where the percent sign is ignored and the
% comments of those files end up being parsed as keys. Touching every language
% of this manual schedules them at the start of the document instead, while the
% percent sign still is a comment character.
\zcsetup{lang=english}
\zcsetup{lang=dutch}
\zcsetup{lang=current}

\newcommand\translation[2]{#1}
\begin{document}
    \selectlanguage{english}
    \DocInput{regulatory-sources.dtx}
\end{document}
%</driver>
% \fi
%
% \subsection{\texorpdfstring{\package{regulatory-sources}}{regulatory-sources}}
% \setcounter{CodelineNo}{0}
%
% A regulatory document cites instruments that are not this document: a statute, a regulation of the
% Union. This module holds the ones a document cites, so that the wording of a citation is a property of
% the document rather than of the sentence it happens to stand in. What it does \emph{not} do is decide the
% wording; that is the next module. This one reads and holds.
%
% \subsubsection{\translation{Prerequisites}{Vereisten}}
%
% \iffalse
%<*package>
% \fi
%    \begin{macrocode}
\RequirePackage{regulatory-struct}
%    \end{macrocode}
%
% \subsubsection{\translation{Sources of a source}{Bronnen van een bron}}
%
% Sources are declared in a \file{bib} file, which is the format such records are already kept in, and
% read by this module rather than by \package{biblatex} or \package{bib2gls}. Neither would be wrong; both
% would put an external program between a document and a handful of records. The test suite of this bundle
% promises a \TeX{} installation and nothing else, the conversion to HTML runs the same documents, and a
% document rarely cites more than a few instruments. So the reading is done here, and the price of that is
% that this module owns a parser.
%
% A parser that is owned has to say what it accepts. This one accepts the shape a \file{bib} file is
% written in when a person writes it, and nothing further:
%
% \begin{labeling}{\quad}
%     \item[] an entry opens with \texttt{@}\meta{type}\texttt{\{}\meta{key}\texttt{,} on a line of its own;
%     \item[] a field is \meta{name}~\texttt{= \{}\meta{value}\texttt{\},} on a line of its own;
%     \item[] the entry closes with \texttt{\}} on a line of its own;
%     \item[] a line that is empty, or that starts with \texttt{\%}, is a comment.
% \end{labeling}
%
% A line that is none of these is an error naming the file and the line, rather than a line that is
% skipped. That is the whole reason a parser may be owned at all: the failure of a reader has to be loud,
% since a record that was not read is a citation that comes out empty and a document that looks finished.
%    \begin{macrocode}
\ExplSyntaxOn
\prop_new:N \g__regulatory_src_required_prop
\prop_new:N \g__regulatory_src_type_prop
\prop_new:N \g__regulatory_src_register_prop
\seq_new:N \g__regulatory_src_keys_seq

\ior_new:N \g__regulatory_src_ior
\str_new:N \l__regulatory_src_key_str
\str_new:N \l__regulatory_src_type_str
\str_new:N \l__regulatory_src_file_str
\str_new:N \l__regulatory_src_query_str
\int_new:N \l__regulatory_src_line_int
\bool_new:N \l__regulatory_src_open_bool
\seq_new:N \l__regulatory_src_parts_seq
\tl_new:N \l__regulatory_src_required_tl
\clist_new:N \l__regulatory_src_required_clist
\tl_new:N \l__regulatory_src_one_tl
\tl_new:N \l__regulatory_src_two_tl
%    \end{macrocode}
%
% \begin{macro}{\newsourcetype}
% Declares an entry type, the fields an entry of it cannot do without, and the register its citations are
% worded in. The required fields are what the wording needs in every register: a Dutch regulation is named
% by its citeertitel, an act of the Union by what kind of act it is and by its number.
%
% The register belongs to the type and not to the document, because it is the legal order of the instrument
% that decides how a citation of it reads. A Dutch contract that cites the \textsc{gdpr} words that
% citation the way the regulation itself does. An entry may override it with a \texttt{register} field.
%    \begin{macrocode}
\NewDocumentCommand \newsourcetype { m m m }
  {
    \prop_gput:Nnn \g__regulatory_src_required_prop { #1 } { #2 }
    \prop_gput:Nnn \g__regulatory_src_register_prop { #1 } { #3 }
  }
\ExplSyntaxOff

\newsourcetype{regeling}{citeertitel}{nl}
\newsourcetype{euhandeling}{soort,nummer}{nl_eu}
%    \end{macrocode}
% \end{macro}
%
% \begin{macro}{\newsource}
% Declares one source in the document itself, which is the route that needs no file at all. It ends in the
% same store as a line read from a \file{bib} file, so nothing downstream can tell the two apart.
%    \begin{macrocode}
\ExplSyntaxOn
\NewDocumentCommand \newsource { m m m }
  {
    \__regulatory_src_begin:nn { #2 } { #1 }
    \keyval_parse:nnn
      { \__regulatory_src_novalue:n } { \__regulatory_src_field:nn } { #3 }
    \__regulatory_src_end:
  }

\cs_new_protected:Npn \__regulatory_src_novalue:n #1
  {
    \msg_error:nnn { regulatory-sources } { no-value } { #1 }
  }
%    \end{macrocode}
% \end{macro}
%
% \subsubsection*{Opening an entry}
%
% \verb|\__regulatory_src_begin:nn|
% Opens an entry of \meta{type} under \meta{key}. An undeclared type and a key that is already taken are
% both errors: the first is a record nothing can word, the second is a record that would quietly replace
% one the document already relies on.
%    \begin{macrocode}
\cs_new_protected:Npn \__regulatory_src_begin:nn #1#2
  {
    \str_set:Nn \l__regulatory_src_type_str { #1 }
    \str_set:Nn \l__regulatory_src_key_str { #2 }
    \prop_if_in:NVF \g__regulatory_src_required_prop \l__regulatory_src_type_str
      {
        \msg_error:nnxx { regulatory-sources } { unknown-type }
          { \l__regulatory_src_type_str } { \l__regulatory_src_key_str }
      }
    \seq_if_in:NVT \g__regulatory_src_keys_seq \l__regulatory_src_key_str
      {
        \msg_error:nnx { regulatory-sources } { duplicate-key }
          { \l__regulatory_src_key_str }
      }
    \prop_new:c { g__regulatory_src_e_ \l__regulatory_src_key_str _prop }
    \bool_set_true:N \l__regulatory_src_open_bool
  }
%    \end{macrocode}
%
% \subsubsection*{Storing a field}
%
% \verb|\__regulatory_src_field:nn|
% Stores one field of the entry that is open. The name of the field is lowered and expanded into the key
% it is stored under, while the value is not expanded at all: a value is data, and a document that writes
% a macro in a bib file gets that macro back rather than what it produced at reading time.
%    \begin{macrocode}
\cs_new_protected:Npn \__regulatory_src_field:nn #1#2
  {
    \bool_if:NTF \l__regulatory_src_open_bool
      {
        \use:x
          {
            \prop_gput:cnn
              { g__regulatory_src_e_ \l__regulatory_src_key_str _prop }
              { \str_lowercase:n { #1 } } { \exp_not:n { #2 } }
          }
      }
      { \msg_error:nnn { regulatory-sources } { stray-field } { #1 } }
  }
%    \end{macrocode}
%
% \subsubsection*{Closing an entry}
%
% \verb|\__regulatory_src_end:|
% Closes the entry and holds it against the fields its type requires. This is the moment the guarantee is
% made: from here on a source may be read without asking whether it is complete, which is what lets the
% commands that word a citation be expandable and silent.
%    \begin{macrocode}
\cs_new_protected:Npn \__regulatory_src_end:
  {
    \prop_get:NVN \g__regulatory_src_required_prop \l__regulatory_src_type_str
      \l__regulatory_src_required_tl
    \quark_if_no_value:NF \l__regulatory_src_required_tl
      {
        \clist_set:NV \l__regulatory_src_required_clist
          \l__regulatory_src_required_tl
        \clist_map_inline:Nn \l__regulatory_src_required_clist
          {
            \prop_if_in:cnF
              { g__regulatory_src_e_ \l__regulatory_src_key_str _prop } { ##1 }
              {
                \msg_error:nnxxx { regulatory-sources } { missing-field }
                  { ##1 } { \l__regulatory_src_key_str }
                  { \l__regulatory_src_type_str }
              }
          }
      }
    \prop_gput:NVV \g__regulatory_src_type_prop
      \l__regulatory_src_key_str \l__regulatory_src_type_str
    \seq_gput_right:NV \g__regulatory_src_keys_seq \l__regulatory_src_key_str
    \bool_set_false:N \l__regulatory_src_open_bool
  }
%    \end{macrocode}
%
% \begin{macro}{\loadsources}
% Reads \meta{file}\texttt{.bib}. The file is looked up the way \TeX{} looks up a file it is given, so it
% is found next to the document or anywhere on the input path; a file that is not there is an error and
% not an empty list of sources.
%    \begin{macrocode}
\NewDocumentCommand \loadsources { m }
  {
    \file_get_full_name:nNTF { #1 .bib } \l__regulatory_src_file_str
      { \__regulatory_src_read:V \l__regulatory_src_file_str }
      { \msg_error:nnn { regulatory-sources } { file-not-found } { #1 .bib } }
  }

\cs_new_protected:Npn \__regulatory_src_read:n #1
  {
    \int_zero:N \l__regulatory_src_line_int
    \bool_set_false:N \l__regulatory_src_open_bool
    \ior_open:Nn \g__regulatory_src_ior { #1 }
    \ior_str_map_inline:Nn \g__regulatory_src_ior
      { \__regulatory_src_line:nn { #1 } { ##1 } }
    \ior_close:N \g__regulatory_src_ior
    \bool_if:NT \l__regulatory_src_open_bool
      {
        \msg_error:nnxx { regulatory-sources } { unclosed }
          { \l__regulatory_src_key_str } { #1 }
      }
  }
\cs_generate_variant:Nn \__regulatory_src_read:n { V }
%    \end{macrocode}
% \end{macro}
%
% \subsubsection*{Reading one line}
%
% \verb|\__regulatory_src_line:nn|
% One line of the file, tried against the four shapes the subset knows. The order is the order they occur
% in: a comment or a blank line first, since a file holds more of those than anything else.
%    \begin{macrocode}
\regex_const:Nn \c__regulatory_src_blank_regex { \A \s* (\%.*)? \Z }
\regex_const:Nn \c__regulatory_src_open_regex
  { \A \s* \@ (\w+) \s* \{ \s* ([^,\s\{\}]+) \s* ,? \s* \Z }
\regex_const:Nn \c__regulatory_src_field_regex
  { \A \s* (\w+) \s* = \s* \{ (.*) \} \s* ,? \s* \Z }
\regex_const:Nn \c__regulatory_src_close_regex { \A \s* \} \s* ,? \s* \Z }

\cs_new_protected:Npn \__regulatory_src_pair:N #1
  {
    \tl_set:Nx \l__regulatory_src_one_tl
      { \seq_item:Nn \l__regulatory_src_parts_seq { 2 } }
    \tl_set:Nx \l__regulatory_src_two_tl
      { \seq_item:Nn \l__regulatory_src_parts_seq { 3 } }
    \exp_args:NVV #1 \l__regulatory_src_one_tl \l__regulatory_src_two_tl
  }

\cs_new_protected:Npn \__regulatory_src_line:nn #1#2
  {
    \int_incr:N \l__regulatory_src_line_int
    \regex_match:NnTF \c__regulatory_src_blank_regex { #2 }
      { }
      {
        \regex_extract_once:NnNTF \c__regulatory_src_open_regex { #2 }
          \l__regulatory_src_parts_seq
          { \__regulatory_src_pair:N \__regulatory_src_begin:nn }
          {
            \regex_extract_once:NnNTF \c__regulatory_src_field_regex { #2 }
              \l__regulatory_src_parts_seq
              { \__regulatory_src_pair:N \__regulatory_src_field:nn }
              {
                \regex_match:NnTF \c__regulatory_src_close_regex { #2 }
                  { \__regulatory_src_end: }
                  {
                    \msg_error:nnxxx { regulatory-sources } { unreadable }
                      { \int_use:N \l__regulatory_src_line_int } { #1 } { #2 }
                  }
              }
          }
      }
  }
%    \end{macrocode}
%
% \begin{macro}{\srcfield}
% \begin{macro}{\srctype}
% Reads one field of a source. Both are expandable and both give nothing for a field that is not there,
% which is safe precisely because a source cannot be declared without the fields its type requires: what
% is absent here is optional by construction.
%    \begin{macrocode}
\cs_new:Npn \srcfield #1#2
  {
    \prop_item:cn { g__regulatory_src_e_ #1 _prop } { #2 }
  }
\cs_new:Npn \srctype #1
  {
    \prop_item:Nn \g__regulatory_src_type_prop { #1 }
  }
%    \end{macrocode}
% \end{macro}
% \end{macro}
%
% \begin{macro}{\srcregister}
% The register a source is worded in: the one its type declares, unless the entry says otherwise.
%    \begin{macrocode}
\cs_new:Npn \srcregister #1
  {
    \tl_if_blank:eTF { \srcfield { #1 } { register } }
      {
        \prop_item:Ne \g__regulatory_src_register_prop
          { \prop_item:Nn \g__regulatory_src_type_prop { #1 } }
      }
      { \srcfield { #1 } { register } }
  }
%    \end{macrocode}
% \end{macro}
%
% \begin{macro}{\ifsourceexists}
% Whether \meta{key} was declared at all. The key is made into a string before it is looked for: the keys
% of a file were read as text, and a sequence compares what it holds character by character, category code
% and all\,---\,so a key typed in the document would never equal the same key read from a file.
%    \begin{macrocode}
\NewDocumentCommand \ifsourceexists { m m m }
  {
    \str_set:Nn \l__regulatory_src_query_str { #1 }
    \seq_if_in:NVTF \g__regulatory_src_keys_seq \l__regulatory_src_query_str
      { #2 } { #3 }
  }
%    \end{macrocode}
% \end{macro}
%
% \subsubsection{\translation{Registers}{Registers}}
%
% A citation is worded in the register of the instrument it points at, and every register has two systems:
% one written out in full and one shortened. Which of the two applies is a property of the place the
% citation stands in\,---\,the Leidraad has it in full in the running text and shortened in a footnote, and
% shortened as well when it stands between brackets in the running text\,---\,so the call decides and the
% document sets the default.
%
% Seven keys are enough to word a citation. Five of them are macros, and each of them is handed the macro
% to put its result in rather than expanding to it: writing an ordinal in words is work that
% \package{fmtcount} does with a command that cannot be expanded, and a citation is assembled before it is
% typeset so that it can be read back.
%
% \begin{labeling}{\texttt{separator}}
%     \item[\texttt{article}] the word before the number of an article.
%     \item[\texttt{number}] the number itself, which is where the colon notation of the civil code lives.
%     \item[\texttt{lid}] the paragraph.
%     \item[\texttt{onderdeel}] the lettered point.
%     \item[\texttt{onder}] the level below that, which the standards call a subonderdeel.
%     \item[\texttt{separator}] what stands between the parts: a comma in full, a space when shortened.
%     \item[\texttt{van}] what stands between the last part and the instrument: \enquote{van} when written
%         out in full, nothing when shortened.
%     \item[\texttt{source}] the instrument itself, named from the fields of its entry.
% \end{labeling}
%    \begin{macrocode}
\ExplSyntaxOff
\newcommand\regulatory@src@ordinal[1]{\ordinalstringnum{#1}}

\newcommand\newsourceregister[3]{%
    \pgfqkeys{/regsrc/#1/voluit}{#2}%
    \pgfqkeys{/regsrc/#1/verkort}{#3}%
}

\newcommand\regsrc@get[3]{\pgfkeysvalueof{/regsrc/#1/#2/#3}}
\newcommand\regsrc@set[4]{\pgfkeysgetvalue{/regsrc/#1/#2/#3}{#4}}
%    \end{macrocode}
%
% The register of the Dutch legislator, which is also the one a Dutch author writes in. In full it is the
% form the Aanwijzingen voor de regelgeving prescribe for regulations themselves\,---\,a model rather than
% a rule for a document that is not a regulation, since those instructions bind ministries and not
% contracts. Shortened it is the form of the Leidraad: no commas, the designation before the number, and
% the colon notation for the civil code.
%    \begin{macrocode}
\newsourceregister{nl}{
    article/.initial=artikel,
    number/.initial=\regulatory@src@number@plain,
    lid/.initial=\regulatory@src@lid@ordinal,
    onderdeel/.initial=\regulatory@src@onderdeel@plain,
    onder/.initial=\regulatory@src@onder@degree,
    separator/.initial={,~},
    van/.initial={van~},
    source/.initial=\regulatory@src@source@full@store
}{
    article/.initial=art.,
    number/.initial=\regulatory@src@number@colon,
    lid/.initial=\regulatory@src@lid@cardinal,
    onderdeel/.initial=\regulatory@src@onderdeel@onder,
    onder/.initial=\regulatory@src@onder@degree,
    separator/.initial={~},
    van/.initial={},
    source/.initial=\regulatory@src@source@short@store
}
%    \end{macrocode}
%
% The register an act of the Union is cited in, which is the one such an act uses for itself: the paragraph
% is a cardinal even in the full form, and the lettered point is a \emph{punt} with the bracket it closes
% with. Counted in the Dutch text of the General Data Protection Regulation: \texttt{punt~a)} thirty-six
% times, \texttt{onderdeel~a)} and \texttt{sub~a)} not once.
%    \begin{macrocode}
\newsourceregister{nl_eu}{
    article/.initial=artikel,
    number/.initial=\regulatory@src@number@plain,
    lid/.initial=\regulatory@src@lid@cardinal,
    onderdeel/.initial=\regulatory@src@onderdeel@punt,
    onder/.initial=\regulatory@src@onder@degree,
    separator/.initial={,~},
    van/.initial={van~},
    source/.initial=\regulatory@src@source@full@store
}{
    article/.initial=art.,
    number/.initial=\regulatory@src@number@plain,
    lid/.initial=\regulatory@src@lid@cardinal,
    onderdeel/.initial=\regulatory@src@onderdeel@punt,
    onder/.initial=\regulatory@src@onder@degree,
    separator/.initial={~},
    van/.initial={},
    source/.initial=\regulatory@src@source@short@store
}
%    \end{macrocode}
%
% The pieces the registers above are built from. A number carries the book of the civil code in the
% shortened form and leaves it to the name of the instrument in the full one, which is what
% \enquote{artikel 96 \dots{} van Boek 6 van het Burgerlijk Wetboek} against \enquote{art. 6:96 BW} comes
% down to.
%    \begin{macrocode}
\newcommand\regulatory@src@number@plain[3]{\def#1{#3}}
\newcommand\regulatory@src@number@colon[3]{%
    \ifthenelse{\equal{#2}{}}{\def#1{#3}}{\def#1{#2:#3}}%
}
\newcommand\regulatory@src@lid@ordinal[2]{%
    \storeordinalstringnum{regsrc}{#2}%
    \expandafter\let\expandafter\regulatory@src@ord\csname @fcs@regsrc\endcsname
    \protected@edef#1{\regulatory@src@ord\space lid}%
}
\newcommand\regulatory@src@lid@cardinal[2]{\def#1{lid~#2}}
\newcommand\regulatory@src@onderdeel@plain[2]{\def#1{onderdeel~#2}}
\newcommand\regulatory@src@onderdeel@onder[2]{\def#1{onder~#2}}
\newcommand\regulatory@src@onderdeel@punt[2]{\def#1{punt~#2)}}
\newcommand\regulatory@src@onder@degree[2]{\def#1{onder~#2\textdegree}}
%    \end{macrocode}
%
% The instrument itself. In full it is named by what it is called, preceded by the article its name takes,
% and the book of a civil code is named here rather than in the number. Shortened it is the abbreviation,
% and a source without one falls back on its name, since a citation that names nothing is worse than a long
% one.
%    \begin{macrocode}
\newcommand\regulatory@src@source@full@store[2]{%
    \protected@edef#1{\regulatory@src@source@full{#2}}%
}
\newcommand\regulatory@src@source@short@store[2]{%
    \protected@edef#1{\regulatory@src@source@short{#2}}%
}

\ExplSyntaxOn
\cs_new:Npn \regulatory@src@lidwoord #1
  {
    \tl_if_blank:eTF { \srcfield { #1 } { lidwoord } }
      { de } { \srcfield { #1 } { lidwoord } }
  }
\cs_new:Npn \regulatory@src@naam #1
  {
    \tl_if_blank:eTF { \srcfield { #1 } { roepnaam } }
      { \srcfield { #1 } { citeertitel } } { \srcfield { #1 } { roepnaam } }
  }
\cs_new:Npn \regulatory@src@source@full #1
  {
    \tl_if_blank:eF { \srcfield { #1 } { boek } }
      { Boek~ \srcfield { #1 } { boek } \c_space_tl van~ }
    \regulatory@src@lidwoord { #1 } ~ \regulatory@src@naam { #1 }
  }
\cs_new:Npn \regulatory@src@source@short #1
  {
    \tl_if_blank:eTF { \srcfield { #1 } { afkorting } }
      { \regulatory@src@naam { #1 } } { \srcfield { #1 } { afkorting } }
  }
\ExplSyntaxOff
\ExplSyntaxOn
%    \end{macrocode}
%
% \subsubsection{\translation{Citing a source}{Een bron aanhalen}}
%
% \begin{macro}{\srcref}
% Cites \meta{key}, either as a whole or down to one of its provisions. The optional argument holds the
% provision\,---\,\texttt{artikel}, \texttt{lid}, \texttt{onderdeel} and \texttt{onder}\,---\,and stands
% before the key, as it does in \cmd{\cite}, so that a bracket in the sentence after the citation is not
% mistaken for it.
%
% The star asks for the system that is not the default of the document, which is what a citation between
% brackets or in a footnote wants. It may also be named outright with \texttt{system}. The default of the
% document is the \option{bronstijl} option, and it is \texttt{voluit}, since that is what the running text
% of a document is.
%
% What it comes to is built up in \cmd{\regulatory@src@last} before it is typeset, so that the wording of a
% citation can be read back rather than only looked at. That is what the test suite asserts on: a citation
% in the wrong register fails nowhere, it simply stands there wrong.
%    \begin{macrocode}
\ExplSyntaxOff
\newcommand\regulatory@src@system{\regulatory@bronstijl}

\define@key{regsrcref}{artikel}{\def\regulatory@src@artikel{#1}}
\define@key{regsrcref}{lid}{\def\regulatory@src@lid{#1}}
\define@key{regsrcref}{onderdeel}{\def\regulatory@src@onderdeel{#1}}
\define@key{regsrcref}{onder}{\def\regulatory@src@onder{#1}}
\define@key{regsrcref}{system}{\def\regulatory@src@thissystem{#1}}

\newcommand\regulatory@src@last{}
\newcommand\regulatory@src@part[1]{%
    \ifx\regulatory@src@parts\@empty\else
        \protected@xdef\regulatory@src@last{\regulatory@src@last\regulatory@src@sep}%
    \fi
    \protected@xdef\regulatory@src@last{\regulatory@src@last#1}%
    \let\regulatory@src@parts\relax
}

\NewDocumentCommand \srcref { s O{} m }{%
    \begingroup
        \def\regulatory@src@artikel{}%
        \def\regulatory@src@lid{}%
        \def\regulatory@src@onderdeel{}%
        \def\regulatory@src@onder{}%
        \IfBooleanTF{#1}{%
            \ifthenelse{\equal{\regulatory@bronstijl}{voluit}}{%
                \def\regulatory@src@thissystem{verkort}%
            }{%
                \def\regulatory@src@thissystem{voluit}%
            }%
        }{%
            \edef\regulatory@src@thissystem{\regulatory@bronstijl}%
        }%
        \setkeys{regsrcref}{#2}%
        \edef\regulatory@src@reg{\srcregister{#3}}%
        \regsrc@set{\regulatory@src@reg}{\regulatory@src@thissystem}{separator}%
            {\regulatory@src@sep}%
        \let\regulatory@src@parts\@empty
        \gdef\regulatory@src@last{}%
        \ifx\regulatory@src@artikel\@empty\else
            \regsrc@set{\regulatory@src@reg}{\regulatory@src@thissystem}{number}{\@@fmt}%
            \@@fmt\regulatory@src@piece{\srcfield{#3}{boek}}{\regulatory@src@artikel}%
            \protected@edef\regulatory@src@piece{%
                \regsrc@get{\regulatory@src@reg}{\regulatory@src@thissystem}{article}%
                \space\regulatory@src@piece}%
            \expandafter\regulatory@src@part\expandafter{\regulatory@src@piece}%
        \fi
        \ifx\regulatory@src@lid\@empty\else
            \regsrc@set{\regulatory@src@reg}{\regulatory@src@thissystem}{lid}{\@@fmt}%
            \@@fmt\regulatory@src@piece{\regulatory@src@lid}%
            \expandafter\regulatory@src@part\expandafter{\regulatory@src@piece}%
        \fi
        \ifx\regulatory@src@onderdeel\@empty\else
            \regsrc@set{\regulatory@src@reg}{\regulatory@src@thissystem}{onderdeel}{\@@fmt}%
            \@@fmt\regulatory@src@piece{\regulatory@src@onderdeel}%
            \expandafter\regulatory@src@part\expandafter{\regulatory@src@piece}%
        \fi
        \ifx\regulatory@src@onder\@empty\else
            \regsrc@set{\regulatory@src@reg}{\regulatory@src@thissystem}{onder}{\@@fmt}%
            \@@fmt\regulatory@src@piece{\regulatory@src@onder}%
            \expandafter\regulatory@src@part\expandafter{\regulatory@src@piece}%
        \fi
        \regsrc@set{\regulatory@src@reg}{\regulatory@src@thissystem}{source}{\@@fmt}%
        \@@fmt\regulatory@src@piece{#3}%
        \ifx\regulatory@src@parts\@empty\else
            \protected@edef\regulatory@src@piece{%
                \regsrc@get{\regulatory@src@reg}{\regulatory@src@thissystem}{van}%
                \regulatory@src@piece}%
        \fi
        \expandafter\regulatory@src@part\expandafter{\regulatory@src@piece}%
    \endgroup
    \regulatory@src@last
}
\ExplSyntaxOn
%    \end{macrocode}
% \end{macro}
%
% \begin{macro}{\regulatory@src@messages}
% What the module says when it stops. Every one of these names the record it was reading, since a bib file
% that is wrong is wrong in one entry and the rest of it is fine.
%    \begin{macrocode}
\msg_new:nnn { regulatory-sources } { unknown-type }
  {
    Source~'#2'~is~of~type~'#1',~which~no~one~declared.~Declare~it~with~
    \iow_char:N\\newsourcetype,~or~use~one~of:~
    \prop_map_function:NN \g__regulatory_src_required_prop
      \__regulatory_src_type_name:nn
  }
\cs_new:Npn \__regulatory_src_type_name:nn #1#2 { ~#1 }

\msg_new:nnn { regulatory-sources } { missing-field }
  {
    Source~'#2'~of~type~'#3'~has~no~'#1'.~A~source~of~that~type~cannot~be~
    written~out~without~it.
  }
\msg_new:nnn { regulatory-sources } { duplicate-key }
  { Source~'#1'~was~already~declared.~The~second~one~would~replace~it. }
\msg_new:nnn { regulatory-sources } { stray-field }
  { Field~'#1'~stands~outside~any~entry. }
\msg_new:nnn { regulatory-sources } { no-value }
  { Field~'#1'~was~given~no~value. }
\msg_new:nnn { regulatory-sources } { file-not-found }
  { There~is~no~file~'#1'~to~read~sources~from. }
\msg_new:nnn { regulatory-sources } { unclosed }
  { Source~'#1'~is~still~open~at~the~end~of~'#2'.~A~closing~brace~is~missing. }
\msg_new:nnn { regulatory-sources } { unreadable }
  {
    Line~#1~of~'#2'~is~none~of~the~four~shapes~this~reader~knows:\\
    ~~#3\\
    An~entry~opens~with~an~at~sign,~the~type~and~the~key;~a~field~reads~
    name~=~value~in~braces;~an~entry~closes~with~a~brace~of~its~own.
  }
\ExplSyntaxOff
%    \end{macrocode}
% \end{macro}
% \iffalse
%</package>
% \fi
%
% \Finale
%
